% c00 —— 卷首（题词、绪言）+ §1 复流形（定义1.1–1.4、习题1.1–1.3），p001–p002
% 原文目次由 main.tex 的 \tableofcontents 自动生成，此处不重排目次。

\vspace{1em}
\begin{quote}
\itshape
The name ``complex group'' formerly advocated by me in allusion to line complexes, as these are
defined by the vanishing of antisymmetric bilinear forms, has become more and more embarrassing
through collision with the word ``complex'' in the connotation of complex number. I therefore
propose to replace it by the corresponding Greek adjective ``symplectic.'' Dickson calls the group
the ``Abelian linear group'' in homage to Abel who first studied it.

\upshape
\raggedleft —— Hermann Weyl \textit{[Wey39, page 165]}
\par
\end{quote}

\vspace{0.5em}
\noindent{\small\color{Sepia}【题词汉译】我以前出于对线汇（line complex，由反对称双线性形式的零点所定义）的呼应，曾把这一群称为「complex group（复群）」；但随着「complex」一词带上复数的含义，这个名称越来越令人尴尬。因此我提议用相应的希腊语形容词「symplectic（辛）」来替代它。Dickson 把该群称为「Abelian linear group（阿贝尔线性群）」，以纪念最早研究它的 Abel。}

\vspace{1em}

复几何起源于代数几何，\footnote{复环面与 K3 曲面中也有非代数的例子，但更显著的例子例如 VII 型曲面等。}但它明确地从代数几何中分离出来，大概是进入 20 世纪以后的事。

另一方面，辛几何起源于 Hamilton 形式的分析力学，而它明确地从分析力学中分离出来，同样也是进入 20 世纪以后的事。

复几何与辛几何自然交汇之处是 Kähler 几何。\footnote{另一个交汇之处是 Hitchin 的广义几何（generalized geometry）。它从一开始就与镜像对称密切相关，但本讲义不涉及。}由 Moishezon 定理，完备的复代数簇若为 Kähler 的，则是射影的。射影簇在复几何、代数几何以及辛几何中都给出了（尽管可能有异议）最重要的一类例子。

「辛（symplectic）」这个词，是 Weyl 作为拉丁语词 complex 在古希腊语中的翻译借用（calque）而造出来的。复几何与辛几何之间存在着超越其词源的神秘关系，这一关系在 20 世纪末被发现。这种关系被称为镜像对称，其发现与超弦理论密切相关。

本讲义从复流形与辛流形的定义讲起，对属于镜像对称的种种现象中最弱的拓扑镜像对称作入门性的解说。\footnote{这是 2020 年度在东京工业大学集中讲义的讲义录。}关于更强的同调镜像对称以及 Strominger–Yau–Zaslow 猜想，例如可参见 [植] 及其参考文献。

\section{复流形}

\begin{definition}
实向量空间 $V$ 的复结构（complex structure）是指 $V$ 的线性自同构 $J\in\End_{\R}(V)$，且满足 $J^2=-\id_V$。
\end{definition}

\begin{exercise}
证明：带有复结构的实向量空间，其作为实向量空间的维数必为偶数。
\end{exercise}

\begin{definition}
对带有复结构的实向量空间 $(V,J_V)$ 与 $(W,J_W)$，从 $(V,J_V)$ 到 $(W,J_W)$ 的态射是指实线性映射 $f\colon V\to W$，且满足
\begin{equation}
    f\circ J_V = J_W\circ f \tag{1.1}
\end{equation}
\end{definition}

\begin{exercise}
证明：带有复结构的实向量空间所构成的范畴，与复向量空间所构成的范畴等价。
\end{exercise}

设 $M$ 为流形，$TM$ 为其切丛。\footnote{本讲义中，除非特别说明，流形（manifold）以及其间的映射等都假定为光滑的（$C^\infty$ 级）。}

\begin{definition}
$\Gamma(\End(TM))$ 中满足 $J^2=-\id_{TM}$ 的元素 $J$，称为 $M$ 的近复结构（almost complex structure）。
\end{definition}

换言之，$M$ 的近复结构就是：在 $M$ 的每一点 $x\in M$ 处切空间 $T_xM$ 的复结构 $J_x$ 所构成的族，且光滑地依赖于 $x$。

\begin{definition}
流形 $M$ 与其上的近复结构 $J_M$ 的组 $(M,J_M)$，称为近复流形（almost complex manifold）。
\end{definition}

% [Tutor Note: 原文 p2 末尾的 演習1.3（对流形 N，切丛 TN 上的正合列 (1.2)）自 p3 起续写，
%  按页边界切分，该题整题归 c01 处理，此处不重复。]
